% Lösungen Einheit 3 von 4 – Der Graphenblick (zusammenbau v0.1)
\einheitenkopf{Einheit 3 von 4 · Der Graphenblick}

\erg{13}{a) $u$ – denn $u' = v$. \quad b) fallend für $x < 1$, steigend für $x > 1$; Tiefpunkt bei $x = 1$ – eine Parabel durch $P(0 | 1)$ mit dem Tiefpunkt $(1 | 0,5)$. \quad c) $F(6) - F(2) = 1 - 3 = -2$ \quad d) $f(2) = 1$ – die Tangente in $(2 | 0)$ hat die Steigung $1$. \quad e) Minimum bei $x = 2$: $g(x) \le 0$ für $x < 2$ und $g(x) > 0$ für $x > 2$, also fällt jede Stammfunktion bis $2$ und steigt danach; die Nullstelle $0$ hat keinen Vorzeichenwechsel und liefert keine Extremstelle. \quad f) wahr – $g$ wechselt bei $-3$ und bei $4$ von minus nach plus (Tiefpunkte), bei $1$ von plus nach minus (Hochpunkt). \quad g) z. B. $a = -1$ und $b = 3$: das Integral ist $f'(3) - f'(-1) = 0 - 0 = 0$.}

13b)

\begin{ksys}[xmin=-2,xmax=4,ymin=-3,ymax=4,klein] \gerade{1}{-1}{f} \funktion{0.5*\x^2-\x+1}{F} \end{ksys}


13d)

\begin{ksys}[xmin=-2,xmax=4,ymin=-2,ymax=4,klein] \parabel{0.5}{1}{-0.5}{F} \tangentean*{0.5*\x^2-\x}{2}{t} \end{ksys}

\erg{14}{a) Jonas hat die Extremstelle von $f$ übernommen; sie ist eine Wendestelle von $F$. Extremstellen von $F$ liegen an den Nullstellen von $f$ mit Vorzeichenwechsel: Hochpunkt bei $x = -2$, Tiefpunkt bei $x = 2$. \quad b) $F' = f$ hat links und rechts der Stelle dasselbe Vorzeichen; $F$ steigt oder fällt auf beiden Seiten weiter – ein Terrassenpunkt, kein Extrempunkt.}
